% Run setup for the beamstrahlung figure (beamstrahlung_fccee.pdf).
% Drop into the lepton-collisions section; \label and citation keys are
% provisional, check them against biblio.bib.

We illustrate the interplay between beamstrahlung, initial-state radiation
and the parton shower with the process $e^-e^+\to\mu^-\mu^+$, generated at
leading order with the dressed-lepton densities of
ref.~\cite{Frixione:2021zdp}.  We consider the two FCC-ee operating points,
the $ZH$ run at $\sqrt{s}=240$~GeV (\texttt{pdlabel}~$=$~\texttt{fcce240ll},
\texttt{ebeam1}~$=$~\texttt{ebeam2}~$=120$~GeV) and the $t\bar t$ run at
$\sqrt{s}=365$~GeV (\texttt{fcce365ll}, $182.5$~GeV per beam), together
with the ILC at $\sqrt{s}=500$~GeV (\texttt{ilc500ll}, $250$~GeV per beam)
as a linear-collider reference.  The beams
are unpolarised and specified with \texttt{lpp1}~$=3$, \texttt{lpp2}~$=-3$;
renormalisation and factorisation scales are dynamical
(\texttt{dynamical\_scale\_choice}~$=-1$, \texttt{scalefact}~$=1$).  The
muons satisfy the default run-card cuts, $p_T>10$~GeV, $|\eta|<2.5$ and
$\Delta R_{\mu\mu}>0.4$, with no cut on $m(\mu^-\mu^+)$.  We generate
$2\times10^4$ unweighted events at each FCC-ee energy and $1.5\times10^4$
for the ILC, obtaining $\sigma=3.358(6)$~pb at $240$~GeV,
$\sigma=1.316(3)$~pb at $365$~GeV and $\sigma=0.3752(6)$~pb at $500$~GeV.

Because a dressed-lepton density convolutes beamstrahlung with ISR, the
momentum fraction carried by the incoming lepton is the product
$x=x_{\rm BS}\,x_{\rm ISR}$.  \MGvii\ samples that splitting event by event
and records $x_{\rm BS1,2}$ in the event file, so that a shower can be told
the beam energy that was actually available at the collision.  Events are
showered with \PYTHIA~8.318~\cite{Bierlich:2022pfr} through the dedicated
main program distributed with the code, which reads those fractions and
degrades each beam to $x_{\rm BS}\,p$ before the event is processed; the
space-like shower then reconstructs $x_{\rm ISR}$ alone.  QED radiation off
leptons is enabled in both showers
(\texttt{SpaceShower:QEDshowerByL}~$=$~\texttt{TimeShower:QEDshowerByL}~$=$
\texttt{on}) together with the lepton densities
(\texttt{PDF:lepton}~$=$~\texttt{on}); multiparton interactions and
hadronisation are switched off, being irrelevant here.

The upper panel of fig.~\ref{fig:beamstrahlung} shows the cumulative
distribution of $E_\gamma^{\rm fwd}$, the summed energy of final-state
photons emitted within $|p_z|/E>0.999$ of either beam, restricted to the
events in which at least one beam radiated.  This is the ISR that the
space-like shower builds.  The dashed curves are obtained by showering the
same events with the beams left at their nominal energy, which is what a
shower unaware of beamstrahlung does.  At the FCC-ee the two treatments are
indistinguishable, because the energy at stake is only
$\sim2\times10^{-5}$ of the beam; at the ILC they separate clearly, the
shower producing on average $8.4$~GeV more forward radiation than it should
against the $9.3$~GeV per event that beamstrahlung had in fact carried away
before the collision, and the fraction of events with no forward photon at
all changing from $26\%$ to $43\%$.  The middle panel shows the underlying
spectra: $5.9\%$ ($6.0\%$) of the FCC-ee beams radiate at $240$ ($365$)~GeV
with a mean relative energy loss of $1.6\times10^{-5}$, against $23.1\%$ and
$7.5\times10^{-3}$ at the ILC.  The lower panel translates this into the
collision energy left after beamstrahlung,
$\sqrt{s'}/\sqrt{s}=\sqrt{x_{\rm BS1}x_{\rm BS2}}$: fewer than $0.2\%$ of
the FCC-ee events fall below $0.999$, whereas $28\%$ of the ILC ones do.
The treatment is therefore essentially inert at a circular collider, as it
should be, and necessary at a linear one.

% Reproduction: vendor/py8main/validation/ holds dump.cc, which showers an
% event file with and without degraded beams and writes the per-event
% quantities, and mg7plot.py, which produces this figure.  Cross sections
% obtained with MadGraph7 0.2.0 at commit 9f9b8cced.

% Cross-check of the shipped spectra against the machine designs.  The
% parameterisations are fits to GuineaPig runs, and the shipped coefficients
% imply, per beam,
%
%   set        <dE>/E    sqrt(<dE^2>)/E     design delta_BS
%   ilc250      0.48%        1.74%          0.65-0.97%
%   ilc500      0.89%        3.14%          1.9%   (4.5% is the 1 TeV row)
%   clic3000    2.69%       11.53%          ~30%
%   fcce240     0.0017%      0.014%         sub-MeV per crossing
%   fcce365     0.0038%      0.024%
%
% For the two ILC points the quoted design value lies between the mean and
% the RMS of the shipped spectrum, and the 500/250 ratio (1.85) matches the
% design one (about 2), so the fits are consistent with the ILC design once
% one is careful about which moment delta_BS denotes.  CLIC3000 is not: even
% its RMS falls a factor three short of the design figure, and that is also
% the only file whose coefficients are written as single-precision literals
% and whose exponential factor is a function of x/z rather than of 1-x/z.
